\documentclass[../main]{subfiles}
\begin{document}
\section{Exposición de Motores de C.A.}
\begin{center}
  \img{images/11/motores.jpg}{Motores de C.A-}
\end{center}
\info{Motores de CA }{Un motor de corriente alterna es una máquina
  que usa electricidad para producir movimiento. A diferencia de la
  corriente continua (como la de las baterías), la corriente alterna
  cambia de dirección constantemente. El motor toma esta corriente
  alterna y la usa para girar un eje o realizar alguna otra tarea mecánica.

  Imagina un ventilador eléctrico que funciona con un enchufe de
  pared. Cuando conectas el ventilador, la electricidad fluye a
  través de él y hace que las aspas giren. Esto es posible gracias a
  un motor de corriente alterna que convierte la electricidad en
  movimiento giratorio.

  En resumen, un motor de corriente alterna es un dispositivo que
  utiliza electricidad que cambia de dirección para crear movimiento,
  como el que encuentras en ventiladores, lavadoras y muchas otras
máquinas eléctricas que usamos a diario.}
\actividad{Exposición de clase especial}
\begin{table}
  \centering
  \begin{tabular}{|c|c|} \hline
    Grupo  & Tema  \\ \hline \hline
    1 &Motor Síncrono \\ \hline
    2& Motor Asíncrono rotor bobinado \\ \hline
    3& Motor Asincrono de jaula de Ardilla \\ \hline
    4 &Motor Universal\\ \hline
    5 &Motor SynRm. \\ \hline
    6& Motor de Fase partida. \\ \hline
  \end{tabular}
\end{table}
\subsection{Recomendaciones de la exposición:}
\begin{itemize}
  \item No leer texto de la presentación, la presentación debe
    contener imágenes y se deben explicar las imágenes y puede
    contener fórmulas gráficos, y tablas, pero no debe tener texto en
    lo posible: sólo títulos.
  \item Pedir el TV una clase antes, y tener el material descargado
    en un pendrive o grabado en una computadora. De forma que no sea
    un problema para la presentación el problema que pueda
    presentarse por no haber disponibilidad del TV o la falta de internet.
  \item Aproximadamente cada alumno debe hablar 5 minutos, la
    presentación no debe ser más extensa que eso por alumno.
  \item En caso de necesitar posponer la fecha de la exposición
    avisar con tiempo al profesor, para organizar una nueva fecha.
\end{itemize}
\end{document}
